% =====================================================================
% Results: Trace-Link Recovery Accuracy and Documentation Fidelity
% IEEEtran. Requires: \usepackage{booktabs}
% =====================================================================
\section{Evaluation: Trace-Link Recovery and Documentation Fidelity}
\label{sec:eval}

We evaluate the pipeline along two complementary axes. The first,
\emph{trace-link recovery accuracy}, asks whether the system localizes the
\emph{correct} source documentation for each API member---i.e., whether the
recovered preprocessed and structured artifacts for member $m$ correspond to the
authoritative documentation of $m$. The second, \emph{documentation fidelity},
asks a conditional question: \emph{given} a correctly recovered trace link, how
faithfully does the structured representation reproduce, place, and retain the
content of the source. Separating the two is essential: a content discrepancy
caused by anchoring to the wrong member (a retrieval failure) is a different
defect from a discrepancy introduced while structuring correctly retrieved text,
and conflating them obscures where the pipeline actually breaks.

\subsection{Corpus and Protocol}
\label{sec:eval-corpus}

We study seven widely used Python libraries---\texttt{numpy} (2.4.0),
\texttt{pandas} (2.3.3), \texttt{requests} (2.32.5), \texttt{scikit-learn}
(1.8.0), \texttt{SQLAlchemy} (2.0.45), \texttt{PyTorch} (2.9.1), and
\texttt{XGBoost} (3.2.0)---each processed from two independent documentation
modalities: the rendered \emph{web} (HTML) documentation and the corresponding
\emph{PDF} manual. Every member is evaluated under both modalities, yielding
$2{,}884$ member--source instances ($1{,}442$ per modality). All reported
numbers are produced under a \emph{full manual census}: every instance was
inspected by a human reviewer rather than sampled. A lightweight, deterministic
triage was used only to \emph{order} and \emph{group} the census so that
near-identical cases could be adjudicated consistently; it makes no contribution
to the reported labels, which are the reviewer's.

\subsection{Trace-Link Recovery Accuracy}
\label{sec:eval-trace}

\subsubsection{Method}
A member's trace link is labelled \emph{correct} when the recovered
documentation unambiguously describes that member, as confirmed against an
authoritative ground truth. Ground truth was the library's installed docstring
and signature where available (the database of record), and, where the static
record was insufficient, the version-pinned official web page for that member.
The reviewer examined every instance the triage marked as correct and confirmed
it, and re-examined every instance the triage flagged as incorrect or uncertain;
this surfaced $29$ triage misclassifications that were corrected before
analysis ($16$ uncertain cases confirmed correct, $7$ presumed-incorrect cases
confirmed correct, and $6$ confirmed incorrect). For correct links we
additionally record a \emph{scope} tag---\emph{full}, \emph{partial}, or
\emph{missing parameters}---capturing whether the entire member documentation was
anchored or only a portion; these flag fidelity caveats and are analyzed in
Section~\ref{sec:eval-fidelity}.

\subsubsection{Results}
Table~\ref{tab:trace-accuracy} reports recovery accuracy per library and
modality. Overall accuracy is $88.5\%$ ($2{,}552/2{,}884$), but the two
modalities diverge sharply: web recovery is near-ceiling at $98.2\%$
($1{,}416/1{,}442$), whereas PDF recovery is $78.8\%$ ($1{,}136/1{,}442$). On
the web modality four of seven libraries are recovered without error, and the
lowest is \texttt{xgboost} at $91.6\%$. The PDF modality is far more variable,
ranging from $48.0\%$ (\texttt{numpy}) to $95.6\%$ (\texttt{SQLAlchemy}).

The failure analysis is unambiguous about \emph{where} the pipeline breaks.
Of the $332$ incorrect links, $327$ ($98.5\%$) are attributable to the
\emph{retrieval/anchoring} stage---the system scraped or anchored to the wrong
span or member---and only $4$ to structuring (an empty or degenerate
extraction), with $1$ residual. Failures are also overwhelmingly concentrated in
the PDF modality ($306$ of $332$), reflecting the weaker structural cues of
linearized PDF text: members that share a name across classes, or whose
anchors fall mid-page, are mis-localized far more often than in HTML, where the
DOM provides explicit member boundaries. This isolates retrieval from
flowed PDF text---not content structuring---as the dominant accuracy bottleneck.

\begin{table}[t]
\centering
\caption{Trace-link recovery accuracy under a full manual census, by library
and documentation modality (correct\,/\,total, percentage). Web recovery is
near-ceiling; PDF recovery is lower and more variable.}
\label{tab:trace-accuracy}
\begin{tabular}{lcc}
\toprule
Library & WEB & PDF \\
\midrule
\texttt{numpy}       & 244/244 (100.0\%) & 117/244 (48.0\%) \\
\texttt{pandas}      & 262/262 (100.0\%) & 218/262 (83.2\%) \\
\texttt{requests}    & 83/87 (95.4\%)    & 65/87 (74.7\%) \\
\texttt{scikit-learn}& 246/246 (100.0\%) & 203/246 (82.5\%) \\
\texttt{SQLAlchemy}  & 153/158 (96.8\%)  & 151/158 (95.6\%) \\
\texttt{PyTorch}     & 276/279 (98.9\%)  & 227/279 (81.4\%) \\
\texttt{XGBoost}     & 152/166 (91.6\%)  & 155/166 (93.4\%) \\
\midrule
\textbf{Overall}     & \textbf{1416/1442 (98.2\%)} & \textbf{1136/1442 (78.8\%)} \\
\bottomrule
\end{tabular}
\end{table}

\subsection{Documentation Fidelity}
\label{sec:eval-fidelity}

Fidelity is evaluated only on the $2{,}552$ correctly recovered trace links, so
that it measures the structuring stage in isolation. A naive verbatim-match
score is uninformative here: a structured field may reproduce its source exactly
in substance yet differ in a handful of characters due to Unicode normalization,
PDF de-hyphenation, smart-quote/dash substitution, or link-placeholder
formatting. We therefore decompose fidelity into three interpretable
sub-constructs and report each as a graded distribution rather than a single
pass/fail rate.

\subsubsection{From manual observation to a deterministic tracer}
\label{sec:eval-fidelity-origin}
The fidelity measurement is grounded in the same manual evaluation that produced
the trace-link census of Section~\ref{sec:eval-trace}. While adjudicating every
recovered link, the reviewer also recorded the \emph{kind} of content
discrepancy each member exhibited, and these qualitative observations---not an
a~priori metric---defined what the measurement needed to capture. Recurring
patterns included: structured signatures that drop a leading
\texttt{class}/\texttt{property}/\texttt{abstract} qualifier or render as a bare
name with empty parentheses; descriptions or parameter/return information
\emph{inferred} from the signature when the source prose documented none; a
member's own purpose populated from an adjacent classmethod's description;
returns descriptions copied from the purpose; usage examples bundled into a
single field, or deposited into a description/\,additional-information field
rather than the examples segment; grouped parameters such as
\texttt{*args,\,**kwargs} either kept under one field or split with a duplicated
shared description; and---pervasively---purely cosmetic divergence from glyph,
hyphenation, and quote normalization, especially in PDF text. These observed
strengths and weaknesses map directly onto three sub-constructs and an
artifact-versus-defect distinction, which we then operationalized as a
\emph{deterministic} tracer (no model is used to score fidelity). Casting the
human-identified phenomena as transparent, reproducible string-grounding rules
lets us report them at corpus scale ($24{,}145$ trace units) while remaining
faithful to what the manual evaluation actually surfaced; spot-checks against the
reviewer's recorded cases were used to set the band thresholds.

\subsubsection{Trace units and sub-constructs}
We score fidelity over the \emph{trace units} that are genuinely shared between
the source code and its API reference: the member \emph{signature}, each
\emph{parameter} (name, type, description), the \emph{returns} (type,
description), the \emph{purpose} (the leading descriptive summary), and
\emph{usage examples}. The three sub-constructs are:
\begin{itemize}
  \item \textbf{Faithfulness (precision):} the fraction of structured trace
  units whose content is grounded in the member's source documentation.
  \item \textbf{Categorization (placement):} of the grounded units, the fraction
  placed in the correct schema field---e.g., a parameter description that is
  grounded in the source but actually copied from the \emph{purpose} prose, or a
  code example deposited in a description field rather than the examples segment,
  is a categorization defect even though its content is faithful.
  \item \textbf{Completeness (recall):} the fraction of \emph{source-documented}
  trace units that are represented in the structured doc.
\end{itemize}

\subsubsection{Measurement}
Operationalizing the manually observed phenomena
(Section~\ref{sec:eval-fidelity-origin}), all comparisons are performed in a
normalized space (Unicode NFKC, de-hyphenation, smart-quote/dash folding,
whitespace collapsing, and unification of URL link placeholders) and entirely
deterministically---no model is used to score fidelity. Each unit is matched against its source by a tiered
procedure: an exact normalized-substring test, then a local fuzzy alignment
(\texttt{rapidfuzz} partial-ratio), then a longest-matching-block fallback. A
unit counts as \emph{faithful} when its best similarity is
$\geq 0.995$ (i.e., differing only by normalization). Categorization compares
each unit against its \emph{expected} source section, identified by structural
headings; placement defects are unit content that grounds in the document at
large but not in its own section. To avoid spurious omissions, completeness is
scoped to trace units rather than arbitrary prose: parameter recall is anchored
on the names declared in the source \emph{signature} (the true code--doc trace
link), and signature, purpose, and returns recall on the presence of those
documented sections. Because the pipeline is prompted to split each example into
its own field, example categorization is evaluated at the level of the examples
\emph{segment} (did example code reach the examples segment at all). Example
\emph{recall} is reported only as a diagnostic: PDF pages concatenate inherited
and sibling members, so example code present on a page frequently belongs to a
different member and cannot be attributed reliably.

To separate genuine structuring defects from measurement artifacts we bin
sub-unit similarity into a \emph{Cosmetic} band ($0.90 \le s < 0.995$),
capturing normalization and glyph drift, and lower bands
(\emph{Paraphrase} $[0.75,0.90)$, \emph{Weak} $[0.55,0.75)$, \emph{Unsupported}
$<0.55$) that constitute real fidelity loss.

\subsubsection{Results: overall}
Table~\ref{tab:fidelity-headline} summarizes the three sub-constructs.
On correctly recovered links the structured documentation is highly faithful
(mean $0.983$ web, $0.955$ PDF), well categorized ($0.977$, $0.973$), and
largely complete ($0.957$, $0.926$). PDF trails web on faithfulness and
completeness but, notably, \emph{not} on categorization: once the correct text
is in hand, the model places it into the right schema field at essentially the
same rate regardless of modality.

\begin{table}[t]
\centering
\caption{Documentation fidelity on correctly recovered trace links. Each cell is
the mean sub-construct score with the percentage of members scoring perfectly
($=1.0$) in parentheses.}
\label{tab:fidelity-headline}
\begin{tabular}{lccc}
\toprule
Source & Faithfulness & Categorization & Completeness \\
\midrule
WEB ($n{=}1413$) & 0.983 (93.0\%) & 0.977 (92.0\%) & 0.957 (87.5\%) \\
PDF ($n{=}1134$) & 0.955 (82.8\%) & 0.973 (90.8\%) & 0.926 (77.6\%) \\
\bottomrule
\end{tabular}
\end{table}

\subsubsection{Results: faithfulness is artifact-dominated}
Table~\ref{tab:fidelity-bands} examines faithfulness at the level of individual
trace units ($13{,}302$ web, $10{,}843$ PDF). Here the headline result is
sharp: $98.6\%$ of web units and $96.9\%$ of PDF units are faithful to
$\geq 0.995$. The largest remaining slice is the \emph{Cosmetic} band
($1.0\%$ web, $2.8\%$ PDF)---normalization and glyph drift, not content error,
and the PDF excess is exactly the expected glyph noise of linearized text.
\emph{Genuine} fidelity defects (Paraphrase and below) account for only
$\approx 0.4\%$ of web units and $\approx 0.3\%$ of PDF units. In other words,
when the structured doc deviates from the source at all, the deviation is almost
always a presentational artifact rather than fabricated or distorted content.

\begin{table}[t]
\centering
\caption{Distribution of structured trace units by similarity to their source
span. The Cosmetic band is normalization/glyph drift (a measurement artifact);
only Paraphrase and below are genuine fidelity defects.}
\label{tab:fidelity-bands}
\begin{tabular}{lrr}
\toprule
Similarity band & WEB (\%) & PDF (\%) \\
\midrule
Faithful ($\geq$0.995)      & 98.6 & 96.9 \\
Cosmetic $[0.90,0.995)$     & 1.0  & 2.8  \\
Paraphrase $[0.75,0.90)$    & 0.2  & 0.1  \\
Weak $[0.55,0.75)$          & 0.2  & 0.1  \\
Unsupported $(<0.55)$       & 0.0  & 0.1  \\
\bottomrule
\end{tabular}
\end{table}

\subsubsection{Results: categorization}
At the member level $92.0\%$ (web) and $90.8\%$ (PDF) of members are perfectly
categorized. The residual defects are of two kinds. \emph{Field placement}
errors are rare and dominated by parameter name/type units
($89/3{,}859$ web; $38/3{,}272$ PDF) and returns descriptions copied from the
purpose prose ($43/573$ web; $41/418$ PDF). \emph{Example misplacement}---code
examples that the source embeds inside descriptive prose and that the model
keeps in a description field rather than promoting to the examples segment---is
the most frequent single categorization issue, affecting $59$ web and $50$ PDF
members; correctly, the model routes example code into the examples segment in
$92.5\%$ (web, $725/784$) and $91.3\%$ (PDF, $527/577$) of members that carry
example code.

\subsubsection{Results: completeness}
Table~\ref{tab:fidelity-completeness} breaks completeness down by trace unit.
Signature and purpose are almost always retained ($\geq 98.6\%$ in both
modalities), parameter recall is high and modality-independent
($93.5\%$ web, $93.4\%$ PDF), and the only material gap is returns recall on PDF
($79.0\%$ vs.\ $92.2\%$ web), where flowed text more often detaches a short
``Returns'' clause from its member. Example recall is reported only as a
diagnostic for the reason given above.

\begin{table}[t]
\centering
\caption{Completeness: percentage of source-documented trace units represented
in the structured doc. Examples are a diagnostic only (PDF pages concatenate
inherited members, inflating apparent omission).}
\label{tab:fidelity-completeness}
\begin{tabular}{lrr}
\toprule
Trace unit & WEB (\%) & PDF (\%) \\
\midrule
Signature        & 99.4 & 99.4 \\
Parameters       & 93.5 & 93.4 \\
Purpose          & 99.4 & 98.6 \\
Returns          & 92.2 & 79.0 \\
\midrule
Examples (diag.) & 73.5 & 57.5 \\
\bottomrule
\end{tabular}
\end{table}

\subsubsection{Results: per-library}
Table~\ref{tab:fidelity-perlib} reports per-library means. Fidelity is high and
stable across libraries; \texttt{scikit-learn} and \texttt{numpy} are strongest,
while \texttt{SQLAlchemy} is the weakest on all three constructs (e.g.,
$0.93/0.89/0.94$ web), consistent with its heavier use of duplicated
table-summary plus standalone member descriptions, which stresses both
faithfulness and placement.

\begin{table}[t]
\centering
\caption{Per-library mean sub-construct scores (F1 Faithfulness / F2
Categorization / F3 Completeness) on correctly recovered trace links.}
\label{tab:fidelity-perlib}
\begin{tabular}{lccc|ccc}
\toprule
& \multicolumn{3}{c|}{WEB} & \multicolumn{3}{c}{PDF} \\
Library & F1 & F2 & F3 & F1 & F2 & F3 \\
\midrule
\texttt{numpy}        & 0.99 & 1.00 & 0.98 & 0.97 & 0.99 & 0.95 \\
\texttt{pandas}       & 0.97 & 1.00 & 0.97 & 0.92 & 0.99 & 0.93 \\
\texttt{requests}     & 0.99 & 0.92 & 0.95 & 0.99 & 0.96 & 0.83 \\
\texttt{scikit-learn} & 0.99 & 1.00 & 0.99 & 0.99 & 1.00 & 0.97 \\
\texttt{SQLAlchemy}   & 0.93 & 0.89 & 0.94 & 0.86 & 0.89 & 0.89 \\
\texttt{PyTorch}      & 1.00 & 0.98 & 0.90 & 0.98 & 0.97 & 0.90 \\
\texttt{XGBoost}      & 0.99 & 0.97 & 0.95 & 1.00 & 0.98 & 0.96 \\
\bottomrule
\end{tabular}
\end{table}

\subsection{Summary}
\label{sec:eval-summary}
Two findings stand out. First, the pipeline's accuracy ceiling is set by
\emph{retrieval}, not structuring: web trace links are recovered with
$98.2\%$ accuracy, almost all of the $332$ failures are PDF anchoring errors
($327/332$ at the retrieval stage), and the dominant lever for improvement is
member localization in linearized PDF text. Second, \emph{conditioned on correct
recovery, the structured output is highly faithful}: $98.6\%$/$96.9\%$ of
web/PDF trace units reproduce their source to within normalization, genuine
content defects are below $0.5\%$, categorization is correct for $\sim 92\%$ of
members, and the only sizable completeness gap is PDF returns recall. The
practical implication is that effort is best directed at robust PDF member
anchoring and at promoting source-embedded example code into the examples
field---not at the faithfulness of the structuring step, which is already at
ceiling.
